%%
%% AgentTelemetry: ASE 2026 Tools & Data Sets Track
%% 4 pages max (including references). Single-blind review.
%% ACM sigconf format with [review] for HotCRP submission.
%%

\documentclass[sigconf,review]{acmart}

%% ---------- packages ----------
\usepackage{listings}
\usepackage{xcolor}
\usepackage{booktabs}
\usepackage{enumitem}
\usepackage{multirow}

%% ---------- listings style ----------
\lstset{
  language=Python,
  basicstyle=\ttfamily\scriptsize,
  keywordstyle=\color{blue!70!black}\bfseries,
  commentstyle=\color{green!50!black},
  stringstyle=\color{red!60!black},
  showstringspaces=false,
  breaklines=true,
  breakatwhitespace=true,
  frame=single,
  framesep=2pt,
  xleftmargin=4pt,
  xrightmargin=2pt,
  numbers=none,
  aboveskip=4pt,
  belowskip=3pt,
  columns=flexible,
}

%% ---------- metadata ----------
\title{AgentTelemetry: An Open-Source Toolkit for Fault Detection in LLM Agent Systems}

\author{Krishna Chaitanya Balusu}
\affiliation{%
  \institution{Independent Researcher}
  \country{United States}
}
\email{krishnabkc15@gmail.com}

\acmConference[ASE '26]{IEEE/ACM 41st International Conference on
Automated Software Engineering}{October 12--16, 2026}{Munich, Germany}
\acmYear{2026}
\copyrightyear{2026}

\begin{abstract}
Autonomous AI agents built on large language models (LLMs) fail in ways
that standard distributed tracing cannot diagnose: reasoning loops,
circular delegation, guardrail bypass, and planning failures all
manifest as opaque \texttt{INTERNAL} spans in OpenTelemetry (OTel).
We present \textsc{AgentTelemetry}, an open-source, pip-installable
Python toolkit (Apache-2.0) that extends OTel with a
nine-kind Domain-Specific Metamodel (DSM) for AI-agent
telemetry---five of the kinds are absent from both the OpenTelemetry
GenAI conventions and OpenInference---enabling typed detector
implementations for 14 agent fault types in a deterministic
fault-injection harness.  The toolkit provides adapters for seven
agent frameworks, four analysis modules, and a real-time circuit
breaker.

In a controlled detector-correctness evaluation under deterministic
fault injection (rate 1.0, mocked LLM clients) across 3{,}528 runs
(14 faults---13 of which align to the MAST multi-agent failure
taxonomy with mapping frozen pre-detector-implementation---$\times$ 6
telemetry conditions $\times$ 7 frameworks $\times$ 6 mocked
foundation models) plus 252 fault-free controls, two AI-agent
baselines (OTel GenAI, OpenInference) and a vanilla-OTel control all
converge to fault detection rate (FDR) 0.429 under typed-only
comparison (detectors must dispatch on a typed enum attribute;
permissive substring-matching is discussed in \S\ref{sec:eval}).
\textsc{AgentTelemetry}'s DSM averages FDR 0.547 on six third-party
adapters (range 0.500--0.571; per-adapter gain +0.07 to +0.14; FPR
rises 0.000$\to$0.083 from one detector on one CrewAI workload).
The qualitative coverage delta is the headline: the DSM is the only
metamodel of the four that can express the eight
coordination/cognitive/policy faults at all (Table~\ref{tab:peradapter}).
On 112 SWE-bench Lite instances a circuit-breaker intervention
improves patch rate by +11~pp (5-seed percentile bootstrap,
approx.\ 95\% CI [+3,+18]) over a no-intervention control with
identical agent/model/tools/prompt/budget; 75\% of failures
(Wilson 95\% CI [66\%, 82\%]) are reasoning loops.  Median
instrumentation overhead 11.7\,\textmu s/span on Apple~M4~Pro and
14.2\,\textmu s on AWS \texttt{c7i.xlarge}.

\smallskip\noindent
\textbf{Tool:} \texttt{pip install agenttelemetry}\\
\textbf{Code:} \url{https://github.com/Krishnachaitanyakc/AgentTelemetry}\\
\textbf{Artifact:} Zenodo DOI \texttt{10.5281/zenodo.PLACEHOLDER}\\
\textbf{Video:} \url{https://youtu.be/PLACEHOLDER}
\end{abstract}

\begin{CCSXML}
<ccs2012>
<concept>
<concept_id>10011007.10011074.10011099</concept_id>
<concept_desc>Software and its engineering~Software verification and validation</concept_desc>
<concept_significance>500</concept_significance>
</concept>
<concept>
<concept_id>10011007.10011074.10011111</concept_id>
<concept_desc>Software and its engineering~Software post-development issues</concept_desc>
<concept_significance>300</concept_significance>
</concept>
</ccs2012>
\end{CCSXML}

\ccsdesc[500]{Software and its engineering~Software verification and validation}
\ccsdesc[300]{Software and its engineering~Software post-development issues}

\keywords{AI agents, observability, OpenTelemetry, fault detection, LLM monitoring, distributed tracing}

\begin{document}

\maketitle

\sloppy

%% ====================================================================
%% 1. INTRODUCTION
%% ====================================================================
\section{Introduction}
\label{sec:intro}

LLM-based autonomous agents---systems that pursue goals through
iterative reasoning, planning, tool use, and delegation---are entering
production deployments in customer support, code review, and incident
response~\cite{cemri2025multiagent}.
Frameworks such as LangChain, CrewAI, AutoGen, and LlamaIndex have
accelerated adoption, but a critical gap remains: \emph{there is no
standardized observability infrastructure for diagnosing agent
failures.}

OpenTelemetry (OTel) is the de-facto standard for distributed tracing.
Its GenAI semantic conventions define operation types for LLM calls,
tool execution, retrieval, and agent lifecycle.  However, five
agent-orchestration phases---planning, reasoning, safety-check
monitoring, inter-agent delegation, and memory management---have no
typed representation in either GenAI or OpenInference, rendering them
as opaque \texttt{INTERNAL} spans.  This semantic gap impedes
automated fault detection for failure modes such as reasoning loops,
circular delegation, guardrail bypass, and planning failures, because
no detector can dispatch on a span kind that the trace never carries.

\textsc{AgentTelemetry} is an open-source Python toolkit for SE
researchers and engineers building, evaluating, or operating LLM-agent
systems who need portable trace data and automated fault triage across
frameworks.  Three contributions:

\begin{enumerate}[leftmargin=*,nosep]
\item \textbf{A Domain-Specific Metamodel (DSM)} of nine agent-specific
  span kinds (five novel) accompanied by six OCL well-formedness
  invariants, derived from systematic analysis of seven production
  frameworks and validated by ablation.
\item \textbf{A modular analysis toolkit} comprising four analysis
  modules (anomaly detection, cost aggregation, decision attribution,
  lexical-overlap grounding screen for triage) and a real-time circuit
  breaker.
\item \textbf{Seven framework adapters} using three instrumentation
  strategies (callback handlers, monkey-patching, manual API),
  requiring as few as 3~lines of code to instrument an agent.
\end{enumerate}

\noindent
The toolkit is pip-installable, OTel-native (all spans export via OTLP
to any compatible backend), and adds 11.7\,\textmu s of median
in-process latency per span---negligible against typical LLM call
latencies of 0.5--10\,s.

%% ====================================================================
%% 2. TOOL ARCHITECTURE
%% ====================================================================
\section{Tool Architecture}
\label{sec:architecture}

\textsc{AgentTelemetry} comprises three modules
(Figure~\ref{fig:arch}): \textbf{Core},
\textbf{Adapters}, and \textbf{Analysis}.

\begin{figure}[t]
\centering
\small
\setlength{\fboxsep}{3pt}
\fbox{\parbox{0.93\columnwidth}{%
\textbf{Core Layer}\\
\textsf{9 Span Kinds} $\cdot$ \textsf{Privacy Filter} $\cdot$
\textsf{Context Propagation} $\cdot$ \textsf{Exporters}\\[3pt]
\textbf{Adapter Layer}\\
\textsf{LangChain} $\cdot$ \textsf{CrewAI} $\cdot$ \textsf{AutoGen}
$\cdot$ \textsf{LlamaIndex} $\cdot$ \textsf{Anthropic} $\cdot$
\textsf{OpenAI} $\cdot$ \textsf{Custom}\\[3pt]
\textbf{Analysis Layer}\\
\textsf{AnomalyDetector} $\cdot$ \textsf{CostAggregator} $\cdot$
\textsf{DecisionAttributor} $\cdot$ \textsf{HallucinationTracer}\\
\textsf{AgentCircuitBreaker} (real-time fault interception)
}}
\caption{Three-layer architecture with OTel-native export.}
\label{fig:arch}
\end{figure}

\subsection{Core: Span Taxonomy}

All nine span kinds are represented as string values in
the \texttt{agenttelemetry.\allowbreak{}span.\allowbreak{}kind}
attribute on standard OTel \texttt{INTERNAL} spans,
ensuring compatibility with OTLP, Jaeger, and any
OTel-compatible backend.
Table~\ref{tab:spans} lists the kinds with their key attributes
and novelty status relative to OTel GenAI.

\begin{table}[t]
\centering
\small
\caption{Nine agent-specific span kinds.
$\star$\,=\,novel (no OTel GenAI equivalent);
ext.\,=\,extends OTel; ovlp.\,=\,overlaps with OTel.}
\label{tab:spans}
\begin{tabular}{@{}llc@{}}
\toprule
\textbf{Span Kind} & \textbf{Key Attributes} & \\
\midrule
\texttt{AGENT}      & \texttt{agent.name, .framework, .role} & ext. \\
\texttt{LLM\_CALL}  & \texttt{llm.model, .tokens, .cost}     & ovlp. \\
\texttt{TOOL\_CALL} & \texttt{tool.name, .input, .status}    & ovlp. \\
\texttt{RETRIEVAL}  & \texttt{retrieval.source, .doc\_count}  & ovlp. \\
\texttt{PLANNING}\,$\star$   & \texttt{planning.strategy, .step\_count} & \\
\texttt{REASONING}\,$\star$  & \texttt{reasoning.chain}               & \\
\texttt{GUARD\_RAIL}\,$\star$& \texttt{guardrail.name, .result}       & \\
\texttt{DELEGATION}\,$\star$ & \texttt{delegation.src/tgt\_agent}     & \\
\texttt{MEMORY}\,$\star$     & \texttt{memory.operation, .key}        & \\
\bottomrule
\end{tabular}
\end{table}

The five novel kinds capture orchestration phases invisible to
OTel GenAI: task decomposition (\texttt{PLANNING}), chain-of-thought
(\texttt{REASONING}), safety-check outcomes (\texttt{GUARD\_\allowbreak{}RAIL}),
inter-agent handoff (\texttt{DELE\-GATION}), and state management
(\texttt{MEMORY}).
An ablation study confirms that every kind is necessary: removing
any one makes at least one of 14~fault types undetectable.

A three-level privacy filter (\texttt{NONE},
\texttt{METADATA\_\allowbreak{}ONLY}, \texttt{FULL}) controls
attribute capture at the span level.  The default
(\texttt{METADATA\_ONLY}) provides full diagnostic signal without
exposing prompt or completion content.

The DSM is also accompanied by six OCL (Object Constraint Language) well-formedness invariants
that the artifact ships as an executable
\texttt{agenttelemetry.ocl} file (loadable into Eclipse OCL or USE).
For example, the \emph{NoCircularDelegation} invariant requires
that no two \texttt{DELEGATION} spans form a 2-cycle on
\texttt{(delegation.src\_agent, delegation.tgt\_agent)};
the others constrain \texttt{LLM\_CALL} cost-attribute presence,
\texttt{GUARD\_RAIL} outcome enum, root-span kind, baggage propagation,
and the kind/subclass discriminator.

\subsection{Adapters: Three Instrumentation Strategies}

Seven framework adapters implement three strategies.
\textbf{Callback handlers} (LangChain, CrewAI, LlamaIndex; 903~LOC)
register event listeners via framework extension APIs.
\textbf{Monkey-patching} (Anthropic SDK, OpenAI SDK, AutoGen; 795~LOC)
wraps key methods at runtime for SDKs without callback systems.
\textbf{Manual API} (Custom; 315~LOC) provides context managers for
explicit instrumentation.
All framework imports are lazy (inside \texttt{\_instrument()}),
avoiding import-time dependencies.


\subsection{Analysis Modules and Circuit Breaker}

Four analysis modules operate on exported span data:
\textbf{CostAggregator} sums token costs by model, agent, and trace;
\textbf{DecisionAttributor} links each \texttt{TOOL\_CALL} to its
parent \texttt{LLM\_CALL} and nearby \texttt{REASONING} spans;
\textbf{AnomalyDetector} detects circular delegation, infinite
retries, cost explosion, and context overflow;
\textbf{HallucinationTracer} provides a lexical-overlap grounding
screen suitable for trace-level triage---it is not a hallucination
classifier and is intended to surface candidate traces for human
review.

The \textbf{AgentCircuitBreaker} is an OTel
\texttt{Span\-Processor} that intercepts fault patterns in real
time.  It dispatches on the typed span kind to check
four policies: reasoning loops ($\geq$\,$N$ identical tool calls),
cost explosion (cumulative cost threshold), context overflow (token
growth rate), and delegation cycles (graph cycle detection).  Each
triggers a configurable callback (log, handler, or exception).  An
equivalent processor over vanilla OTel spans is possible but requires
heuristic parsing of unstructured span names and free-form attributes;
typed dispatch makes the policy code one line per fault.

\paragraph*{Requirements.} Python 3.10+; \texttt{opentelemetry-sdk}
$\geq$1.27; the seven framework adapters lazily import their target
SDK when \texttt{instrument()} is called, so installs do not pull in
all seven by default.  An \texttt{extras\_require} group exists per
framework (e.g.\ \texttt{pip install "agenttelemetry[langchain]"}).
The artifact ships a Dockerfile and a \texttt{docker-compose.yml} that
spins up an OTel collector and Jaeger backend reproducing all numbers
in \S\ref{sec:eval}.

%% ====================================================================
%% 3. USAGE
%% ====================================================================
\section{Usage}
\label{sec:usage}

After \texttt{pip install agenttelemetry}, basic instrumentation is
three lines (Listing~\ref{lst:usage}); framework adapters auto-hook
on a single \texttt{adapter.instrument()} call; spans export via OTLP
to Jaeger, Grafana Tempo, Datadog, or any OTel-compatible backend; a
JSON-Lines exporter supports offline analysis.

\begin{lstlisting}[language=Python,basicstyle=\ttfamily\tiny,frame=single,framesep=2pt,xleftmargin=2pt,caption={Three-line install + circuit-breaker.},label={lst:usage}]
from agenttelemetry import AgentTracer, AgentCircuitBreaker
tr = AgentTracer("research-agent")
tr.add_span_processor(AgentCircuitBreaker(max_reasoning_loops=3,
    max_cost=0.50, on_trigger=lambda f: print(f)))  # demo callback
with tr.start_agent_span("researcher"):
    with tr.start_llm_span("gpt-4o"): r = llm.chat(prompt)
    with tr.start_tool_span("search"): d = search_tool(query)
\end{lstlisting}

%% ====================================================================
%% 4. EVALUATION
%% ====================================================================
\section{Evaluation}
\label{sec:eval}

We evaluate \textsc{AgentTelemetry} along four dimensions: fault
detection effectiveness, span kind necessity, runtime overhead, and
real-world applicability.

\subsection{Fault Detection Rate}

We define 14 fault types spanning all nine span kinds (e.g., reasoning
loop, circular delegation, guardrail bypass, planning failure, cost
explosion, hallucination); 13 of the 14 map onto modes of the MAST
multi-agent failure taxonomy~\cite{cemri2025multiagent} (the 14th,
\texttt{COST\_EXPLOSION}, was added independently from public
cost-runaway postmortems for autonomous-agent stacks).  The
taxonomy was finalized \emph{before} we wrote the detectors, and the
9-kind DSM was finalized \emph{before} the taxonomy was frozen; the
artifact (\S\ref{sec:concl}) ships a \texttt{mast\_mapping.md} with
git-commit timestamps establishing this order (DSM
$\to$ taxonomy $\to$ detector implementation).
Each fault is injected at rate~1.0 into mock LLM clients.  FDR here
is operationalized as: under each metamodel, does a typed-enum-dispatch
detector exist that fires on the injected fault?  Baselines fail not
because the underlying signal is absent from the trace's free-form
data but because no typed attribute exists for the detector to
dispatch on; permissive substring-matching scoring is discussed below.

Across 3{,}528 fault-injection runs and 252 fault-free controls, the
two AI-agent baselines (OTel GenAI, OpenInference) and a vanilla-OTel
control all converge to FDR\,=\,0.429 under
typed-only comparison.  All three detect the same six tool/LLM-attribute
faults: \texttt{wrong\_tool}, \texttt{tool\_failure}, \texttt{timeout},
\texttt{infinite\_loop}, \texttt{context\_overflow}, \texttt{cost\_explosion}.
None expresses the eight coordination/cognitive/policy faults; in
particular, OpenInference's typed \texttt{GUARDRAIL} kind does not
yield \texttt{guardrail\_bypass} detection because the spec defines no
constrained outcome attribute, and its typed \texttt{AGENT} kind does
not yield \texttt{agent\_misroute} detection because no typed
inter-agent edge is captured.  \textsc{AgentTelemetry}'s DSM is the
only metamodel of the four that \emph{can} express those eight.  It
reaches a constructive upper bound of FDR\,=\,1.000 on the
conformance-complete reference adapter (\texttt{custom\_agent}) and
averages 0.547 across the six third-party adapters (Tables~\ref{tab:fdr},~\ref{tab:peradapter}).  A permissive scoring rule allowing substring-matching on \texttt{span.name} would lift the baselines but at higher per-framework FPR; we use typed-enum-dispatch for the apples-to-apples metamodel comparison.
False-positive rate is 0.000 for the three baselines and 0.083 for
\textsc{AgentTelemetry}, traced to the \texttt{infinite\_loop} detector
misfiring on a CrewAI workload that legitimately repeats a tool call
across writer turns---a thresholding artifact, addressable with a
state-aware OCL constraint (future work).

\begin{table}[t]
\centering
\small
\caption{Aggregate FDR / FPR by telemetry condition over 3{,}528
fault-injection runs and 252 fault-free controls.  The headline DSM
row is the third-party-adapter mean (n=6, the honest generalization).
The conformance-complete reference adapter achieves 1.000 by
construction (we built it to satisfy the OCL invariants); we report
it in the ablation context (\S\ref{sec:ablation}) rather than as a
headline number.}
\label{tab:fdr}
\begin{tabular}{@{}lcc@{}}
\toprule
\textbf{Condition} & \textbf{FDR} & \textbf{FPR} \\
\midrule
No telemetry              & 0.000 & 0.000 \\
Vanilla OTel              & 0.429 (6/14) & 0.000 \\
OTel GenAI (current spec) & 0.429 (6/14) & 0.000 \\
OpenInference             & 0.429 (6/14) & 0.000 \\
\textbf{DSM, third-party adapters (n=6 mean)} & \textbf{0.547} & \textbf{0.083} \\
\bottomrule
\end{tabular}
\end{table}

\begin{table}[t]
\centering
\scriptsize
\caption{Which fault types each metamodel can \emph{express} (\checkmark) under typed-only dispatch.  The DSM uniquely covers the eight coordination/cognitive/policy faults below the rule.  Per-adapter FDR for AgentTelemetry: \texttt{custom} 1.000 (by construction), \texttt{anthropic\_sdk}/\texttt{autogen}/\texttt{crewai}/\texttt{openai\_sdk} 0.571, \texttt{langchain}/\texttt{llamaindex} 0.500.}
\label{tab:peradapter}
\begin{tabular}{@{}lcccc@{\quad}lcccc@{}}
\toprule
\textbf{Fault} & \textbf{V} & \textbf{G} & \textbf{O} & \textbf{D} & \textbf{Fault} & \textbf{V} & \textbf{G} & \textbf{O} & \textbf{D} \\
\midrule
wrong\_tool       & \checkmark & \checkmark & \checkmark & \checkmark & circular\_deleg.  & -- & -- & -- & \checkmark \\
tool\_failure     & \checkmark & \checkmark & \checkmark & \checkmark & agent\_misroute   & -- & -- & -- & \checkmark \\
timeout           & \checkmark & \checkmark & \checkmark & \checkmark & planning\_fail.   & -- & -- & -- & \checkmark \\
infinite\_loop    & \checkmark & \checkmark & \checkmark & \checkmark & reasoning\_loop   & -- & -- & -- & \checkmark \\
context\_overfl.  & \checkmark & \checkmark & \checkmark & \checkmark & guardrail\_byp.   & -- & -- & -- & \checkmark \\
cost\_explosion   & \checkmark & \checkmark & \checkmark & \checkmark & stale\_retr.      & -- & -- & -- & \checkmark \\
                  &            &            &            &            & memory\_corr.     & -- & -- & -- & \checkmark \\
                  &            &            &            &            & hallucination     & -- & -- & -- & \checkmark \\
\midrule
\multicolumn{10}{@{}p{0.95\columnwidth}@{}}{\textbf{Coverage:} V-OTel 6/14, GenAI 6/14, OpenInference 6/14, \textbf{DSM 14/14}.  V=Vanilla OTel, G=GenAI, O=OpenInference, D=DSM.} \\
\bottomrule
\end{tabular}
\end{table}

\subsection{Ablation Study}
\label{sec:ablation}

For each of the nine span kinds, stripping all spans of that kind from
the exported trace makes at least one fault undetectable.  AGENT loses
\texttt{agent\_misroute}; LLM\_CALL loses \texttt{cost\_explosion},
\texttt{context\_overflow}, \texttt{hallucination}, \texttt{infinite\_loop};
TOOL\_CALL loses \texttt{wrong\_tool}, \texttt{tool\_failure},
\texttt{infinite\_loop}; RETRIEVAL loses \texttt{stale\_retrieval},
\texttt{hallucination}; PLANNING/REASONING/GUARD\_RAIL/MEMORY each lose
their namesake fault; DELEGATION loses \texttt{circular\_delegation}
and \texttt{agent\_misroute}.  No kind can be removed without losing
at least one detection capability.

\subsection{SWE-bench Case Study}

On 112 SWE-bench Lite instances (37\% of the benchmark, all 12
repositories), we instrument a ReAct coding agent (GPT-4o-mini, 5
tools, 8 iterations max, \$2.53 total API cost).  The agent produces
3,060 spans across all nine kinds.

\textsc{AgentTelemetry} characterizes reasoning loops as the dominant
failure mode: 75\% of 84 failed instances (Wilson 95\% CI [66\%, 82\%])
repeatedly called the same tool with similar queries without
converging.  We define a reasoning loop as $\geq$3 \texttt{TOOL\_CALL}
spans whose \texttt{tool.name} and a 3-token-shingle Jaccard
similarity~$\geq$0.7 over the input arguments are identical; manual
audit of 20 randomly-selected loop-flagged traces by the author gave
18 true loops (precision 0.90, single-author audit; inter-rater
reliability not measured---the 75\% reasoning-loop attribution
inherits this measurement-validity gap; two false positives were
legitimate refinement of an
edit-file argument).  This pattern is invisible to vanilla OTel, which
cannot distinguish a reasoning step from any other \texttt{INTERNAL}
span.

A telemetry-guided intervention---the circuit breaker injects a
strategy-change prompt (verbatim: \emph{``You have repeated the same
tool call three times without progress.  Stop, reconsider the goal,
and try a different tool or argument shape''}) after a flagged
loop---improves the patch rate by +11~percentage points (5-seed
percentile bootstrap of per-seed differences, approximate 95\% CI
[+3, +18]; baseline 25.0\% (SD across 5 seeds: 3.4\%), intervention
36.0\% (SD: 3.1\%); per-seed pairs reported in the artifact).  Each
seed runs all 112 instances with identical agent configuration
(GPT-4o-mini, 5 tools, 8-iteration max, identical
system prompt, fixed token budget); the only delta is whether the
circuit breaker's \texttt{on\_trigger} callback is enabled.  Per-seed
logs and the exact prompt are in
\texttt{experiments/swebench\_closed\_loop.py}.
A sham-prompt arm (non-loop-related reflection prompt at the same
trigger points) is in the artifact roadmap and will isolate the
detection contribution from the reflection-prompt contribution.

\subsection{Overhead}

Measured over 90{,}000 iterations (10{,}000 per kind) on two reference
machines: an Apple~M4~Pro laptop and a 4~vCPU AWS \texttt{c7i.xlarge}
cloud VM.  Median span creation is 11.7\,\textmu s (M4) /
14.2\,\textmu s (\texttt{c7i}); p99 28.2\,/\,35.1\,\textmu s; memory
$\sim$3.1\,KB per span; sustained throughput 58--66\,k\,/\,42--48\,k
spans/s.  Given 1--10 spans/s and LLM latencies of 0.5--10\,s, the
per-span overhead adds $<$0.006\% to end-to-end execution.

%% ====================================================================
%% 5. RELATED WORK
%% ====================================================================
\section{Related Work}
\label{sec:related}

Table~\ref{tab:comparison} positions \textsc{AgentTelemetry} against
representative tools across seven dimensions critical to production
agent observability.

\begin{table}[t]
\centering
\scriptsize
\caption{Comparison with existing tools.
\checkmark\,=\,supported; $\circ$\,=\,partial; --\,=\,absent.
\textbf{Span kinds} = typed agent-orchestration kinds beyond
LLM/tool/retrieval.  \textbf{Privacy} = declarative per-span
attribute-capture levels with default redaction.
\textbf{Mng. SaaS} = managed service with SLA.  AgentTelemetry is
research-grade software; LangSmith and Datadog provide commercial
support and substantially larger production user bases.}
\label{tab:comparison}
\begin{tabular}{@{}lcccccccc@{}}
\toprule
\textbf{Tool} & \rotatebox{70}{\textbf{Multi-fw}} &
\rotatebox{70}{\textbf{Span kinds}} &
\rotatebox{70}{\textbf{Multi-agent}} &
\rotatebox{70}{\textbf{Privacy}} &
\rotatebox{70}{\textbf{Fault det.}} &
\rotatebox{70}{\textbf{Open-src}} &
\rotatebox{70}{\textbf{OTel-native}} &
\rotatebox{70}{\textbf{Mng. SaaS}} \\
\midrule
LangSmith~\cite{langsmith}     & --          & $\circ$ & -- & -- & -- & -- & -- & \checkmark \\
Langfuse~\cite{langfuse}       & $\circ$     & --      & -- & -- & -- & \checkmark & \checkmark & \checkmark \\
OpenLLMetry~\cite{openllmetry} & \checkmark  & --      & -- & -- & -- & \checkmark & \checkmark & -- \\
Datadog LLM~\cite{datadog_llm} & \checkmark  & --      & -- & -- & -- & -- & $\circ$ & \checkmark \\
OTel GenAI~\cite{otelgenai}    & \checkmark  & $\circ$ & $\circ$ & -- & -- & \checkmark & \checkmark & -- \\
OpenInference~\cite{openinference} & \checkmark & $\circ$ & $\circ$ & -- & -- & \checkmark & \checkmark & \checkmark \\
AGDebugger~\cite{agdebugger}   & --          & $\circ$ & \checkmark & -- & $\circ$ & \checkmark & -- & -- \\
\midrule
\textbf{AgentTelemetry}        & \checkmark  & \checkmark & $\circ$ & \checkmark & \checkmark & \checkmark & \checkmark & -- \\
\bottomrule
\end{tabular}
\end{table}

\textbf{Failure taxonomies.}
The MAST taxonomy of Cemri et al.~\cite{cemri2025multiagent} identifies
14~multi-agent failure modes from 200+ annotated traces; 13/14 of our
fault types map to MAST modes (artifact: \texttt{mast\_mapping.md}).
These taxonomies characterize \emph{what} goes wrong but provide no
runtime detection infrastructure.

\textbf{Debugging tools.}
AGDebugger~\cite{agdebugger} offers interactive visual debugging for
multi-agent conversations.  It operates on framework-specific logs;
\textsc{AgentTelemetry} provides standardized, OTel-native trace data
that such tools could consume.

\textbf{LLM observability platforms.}
LangSmith~\cite{langsmith} provides deep LangChain integration but
no portability to other frameworks.
Langfuse~\cite{langfuse} and
OpenLLMetry~\cite{openllmetry} provide OTel-compatible tracing but
define no agent-specific span kinds---a planning step and a
summarization call are indistinguishable.
Datadog LLM Observability~\cite{datadog_llm} adds LLM metrics to a
commercial stack but cannot detect agent orchestration faults.

\textbf{Industry typed-span conventions.}
The OTel GenAI conventions~\cite{otelgenai} define typed operations
covering LLM calls, retrieval, tools, agent creation, and agent
invocation.  OpenInference~\cite{openinference} defines ten typed span
kinds including \texttt{AGENT} and \texttt{GUARDRAIL}.  Both lift
agent and tool to first-class types but neither exposes typed
sub-operations for delegation topology, planning, reasoning, or
memory, nor a constrained guardrail outcome.  \textsc{AgentTelemetry}
is a strict typed-superset and coexists on the same OTLP export
pipeline.

%% ====================================================================
%% 6. CONCLUSION
%% ====================================================================
\section{Conclusion}
\label{sec:concl}

\textsc{AgentTelemetry} extends OpenTelemetry with a nine-kind DSM,
six OCL invariants, seven framework adapters, four analysis modules,
and a real-time circuit breaker.  Under typed-only comparison the
DSM averages FDR 0.547 on third-party adapters (vs.\ 0.429 for vanilla
OTel / OTel GenAI / OpenInference) and reaches 1.000 on a
conformance-complete reference adapter; on SWE-bench Lite a
circuit-breaker intervention improves patch rate by +11~pp
(approx.\ 95\% CI [+3,+18], 5 seeds).  Median per-span overhead is
$<$15\,\textmu s on commodity x86 cloud.

\paragraph*{Limitations.}
Mocked LLMs with rate-1 fault injection (not real production faults);
single-agent SWE-bench (5 seeds, lower bound +3pp); no sham-prompt
control isolating detection from prompt-injection; single-author
loop-detector audit (n=20); third-party adapter coverage uneven by
framework, so the third-party gain reflects today's adapter realization,
not the metamodel ceiling; FPR gap (0.083 vs.\ 0.000) localizes to one
CrewAI thresholding artifact, not a metamodel limitation.

The toolkit is on PyPI (\texttt{pip install agenttelemetry},
Apache-2.0); code at
\url{https://github.com/Krishnachaitanyakc/AgentTelemetry}; artifact
at Zenodo DOI \texttt{10.5281/zenodo.PLACEHOLDER}.  Future work:
state-aware OCL to remove the CrewAI false positive; sham-prompt
control on the SWE-bench intervention; permissive-baseline study;
adaptive sampling for long-running agents.

%% ====================================================================
%% BIBLIOGRAPHY
%% ====================================================================
\begin{thebibliography}{10}

\bibitem{cemri2025multiagent}
M.~Cemri, M.~Z. Pan, S.~Yang, et al.
\newblock Why Do Multi-Agent {LLM} Systems Fail?
\newblock {\em arXiv:2503.13657}, 2025.

\bibitem{agdebugger}
W.~Epperson, G.~Bansal, V.~Dibia, et al.
\newblock Interactive Debugging and Steering of Multi-Agent {AI}
Systems.
\newblock In {\em Proc.\ CHI}, 2025.

\bibitem{langsmith}
LangChain.
\newblock {LangSmith}: LLM Application Observability Platform.
\newblock \url{https://smith.langchain.com}, 2024.

\bibitem{langfuse}
{Langfuse}: Open-Source {LLM} Observability.
\newblock \url{https://langfuse.com}, 2024.

\bibitem{openllmetry}
Traceloop.
\newblock {OpenLLMetry}: OpenTelemetry-based observability for {LLM}
applications.
\newblock \url{https://github.com/traceloop/openllmetry}, 2024.

\bibitem{datadog_llm}
Datadog.
\newblock {LLM Observability}.
\newblock \url{https://www.datadoghq.com/product/llm-observability/}, 2024.

\bibitem{otelgenai}
{OpenTelemetry Contributors}.
\newblock {GenAI} Semantic Conventions.
\newblock \url{https://opentelemetry.io/docs/specs/semconv/gen-ai/}, 2024.

\bibitem{openinference}
{Arize AI}.
\newblock {OpenInference}: Semantic Conventions for {AI} System
Observability.
\newblock \url{https://github.com/Arize-ai/openinference}, 2024.

\end{thebibliography}

\end{document}
